\section{The target: ECDSA, secp256k1, and one reversible point addition}\label{sec:target}

\subsection{What a quantum computer has to compute}

ECDSA on secp256k1 signs with a secret scalar $k\in[1,r)$ whose public key is $Q=[k]P$, where $P$ is the standard generator of the curve
\[
E:\ y^2=x^3+7 \quad\text{over } \F_p,\qquad p=2^{256}-2^{32}-977,\qquad r=\#E(\F_p)\ \text{prime}.
\]
Verifying a signature needs only $Q$; forging one needs $k$; and recovering $k$ from $(P,Q)$ is the elliptic-curve discrete logarithm problem (ECDLP). Classically the best attacks are generic $\Theta(\sqrt r)\approx2^{128}$ group operations. Shor's algorithm \cite{shor97,proos-zalka03} solves it in polynomial time by phase estimation on the two-dimensional ``probe'' function
\begin{equation}\label{eq:probe}
f(u,v)=[u]P+[v]Q=[u+kv]P,\qquad f(u-k,\,v+1)=f(u,v),
\end{equation}
whose period lattice is generated by $(-k,1)$ and $(r,0)$: prepare two uniform superpositions over $u,v\in[0,2^m)$ with $2^m\gtrsim r$, compute $f$ coherently into a point register, Fourier-transform both scalar registers, and measure a pair $(s,t)$ that satisfies $t\equiv ks\pmod r$ with constant probability, from which $k=ts^{-1}\bmod r$ and the candidate is verified classically by $[k]P\overset{?}{=}Q$ \cite{ekera19}.

\begin{figure}[H]
\centering
\resizebox{\textwidth}{!}{%
\begin{quantikz}[column sep=6pt,row sep=5pt,font=\small]
\lstick{$\ket{u_0}$} & \qwbundle{16} & \gate{H^{\otimes16}} & \gate{\mathrm{lookup}} & \ctrl{6} & \gate{\mathrm{unlookup}} &          &         &          & \gate[3]{\mathcal F^{-1}} & \meter{} \\
\lstick{$\vdots$}     &  &                      &  &  &  & \ \cdots\ &         &          &                           &          \\
\lstick{$\ket{u_{15}}$} & \qwbundle{16} & \gate{H^{\otimes16}} &  &  &  &          & \ctrl{4}&          &                           & \meter{} \\
\lstick{$\ket{v_0}$} & \qwbundle{16} & \gate{H^{\otimes16}} &  &  &  &          &         & \ctrl{3} & \gate[3]{\mathcal F^{-1}} & \meter{} \\
\lstick{$\vdots$}     &  &                      &  &  &  &          &         &          &                           &          \\
\lstick{$\ket{v_{15}}$} & \qwbundle{16} & \gate{H^{\otimes16}} &  &  &  &          &         &          &                           & \meter{} \\
\lstick{$R:\ \ket{S}$} & \qwbundle{2\times256} &                   &  & \gate{+\,T_0[u_0]} & & \ \cdots\ & \gate{+\,T_{15}[u_{15}]} & \gate{+\,T'_{0}[v_0]} & \ \cdots &
\end{quantikz}}
\caption{Shor's ECDLP circuit with the scalars split into $16$-bit windows ($n=256$, $w=16$): each window $z\in\{u_j,v_j\}$ is a coherent address into a classical table $T_j[i]=[i\,2^{16j}]P$ (or $[i\,2^{16j}]Q$); a QROM lookup materialises the selected point, the accumulator $R$ is updated in place by one \emph{mixed} point addition, and the lookup is uncomputed. The scalar registers pass through unchanged and are read out by (semiclassical) inverse Fourier transforms. The accumulator starts at a fixed point $S$, never at $\mathcal O$, so that the affine formulas stay valid \cite{schrottenloher26,haener26}. Redrawn after Supplemental Note~D of \cite{ecdsafail26}.}
\label{fig:shor-ecdlp}
\end{figure}

\paragraph{Windowing and the cost formula.}
With $w$-bit windows there are $\lceil n/w\rceil=16$ windows per scalar, i.e.\ $32$ conceptual point additions; the first window is replaced by a direct table lookup of the $2^{2w}$ possible starting points (or $2^{18}$ entries in IonQ's version \cite{haener26}) and the last three windows are dropped and recovered by a classical search over $2^{48}$ candidates, which leaves $2n/w-4=28$ executed point additions \cite{babbush26,haener26,ecdsafail26}. Each addition also pays three table lookups of $2^{16}$ entries (the coordinates $x_A,y_A$ and the derived value $3x_A$), which both IonQ and \ecdsafail{} price at $3\cdot2^{16}$ Toffolis. This gives the proxy that \ecdsafail{} uses to convert its per-addition numbers into whole-attack numbers (its eq.~``full-attack''):
\begin{equation}\label{eq:full-attack}
\text{Toffolis}\approx 28\,(T+3\cdot2^{16}),\qquad \text{logical qubits}\approx Q+16,
\end{equation}
and, if one point addition fails on a fraction $e$ of inputs, a union bound $1-28e$ on the success probability. The last three windows can be dropped because the classical post-processing verifies each candidate against the curve; Eker{\aa}'s analysis \cite{ekera19} shows a few extra qubits per scalar register restore the success probability that the $r\nmid2^m$ rounding would otherwise lose. The Fourier transforms contribute nothing to the Toffoli count and can be done semiclassically on one recycled qubit per register \cite{griffiths-niu96}; every published estimate therefore reports only the point arithmetic.

\subsection{The affine group law and the in-place mixed addition}

For $R=(x_R,y_R)$ and $A=(x_A,y_A)$ with $x_R\neq x_A$,
\begin{equation}\label{eq:grouplaw}
\lambda=\frac{y_R-y_A}{x_R-x_A},\qquad x_{R'}=\lambda^2-x_R-x_A,\qquad y_{R'}=\lambda(x_A-x_{R'})-y_A .
\end{equation}
Three subtractions, one squaring, one multiplication and one \emph{division}. The division is the whole problem: the divisor $x_R-x_A$ is a quantum value, so the inverse has to be computed by a reversible circuit on a superposition and then uncomputed. \ecdsafail{}'s best circuits implement the group law as the seven register updates of \cref{tab:seven-steps}, which overwrite the $X$ and $Y$ registers in place and never allocate separate output coordinates (Supplemental Note~D, \cite{ecdsafail26}).

\begin{table}[H]
\centering\small
\begin{tabular}{@{}clll@{}}
\toprule
step & update & $X$ register afterwards & $Y$ register afterwards\\
\midrule
1 & $X\gets x_R-x_A;\ \ Y\gets y_R-y_A$ & $d_x=x_R-x_A$ & $d_y=y_R-y_A$\\
2 & $Y\gets Y\cdot X^{-1}$ & $d_x$ & $\lambda$\\
3 & $X\gets X+3x_A$ & $x_R+2x_A$ & $\lambda$\\
4 & $X\gets X-Y^2$ & $x_A-x_{R'}$ & $\lambda$\\
5 & $Y\gets Y\cdot X$ & $x_A-x_{R'}$ & $\lambda(x_A-x_{R'})$\\
6 & $Y\gets Y-y_A$ & $x_A-x_{R'}$ & $y_{R'}$\\
7 & $X\gets x_A-X$ & $x_{R'}$ & $y_{R'}$\\
\bottomrule
\end{tabular}
\caption{In-place mixed addition $R'=R+A$ in seven reversible register updates (Supplemental Note~D of \cite{ecdsafail26}, after \cite{schrottenloher26}). Steps 2 and 5 are the two multiplicative steps; after steps 3--4, $X=(x_R-x_A)+3x_A-\lambda^2=x_A-x_{R'}$. Step 5 is \emph{exactly} the inverse of the division unitary of step 2 applied to a different operand: $U_{\rm div}:\ket{X}\ket{Y}\mapsto\ket{X}\ket{YX^{-1}}$, so $U_{\rm div}^{-1}:\ket{X}\ket{Y}\mapsto\ket{X}\ket{YX}$.}
\label{tab:seven-steps}
\end{table}

\begin{warnbox}[title={The generic-case domain}]
The seven steps are correct only when $x_R\neq x_A$ (step 2 divides by $d_x$) and $x_A\neq x_{R'}$ (step 5 multiplies by $x_A-x_{R'}$; in the register-shared low-$Q$ circuit this quantity is also inverted). The excluded inputs $R\in\{A,-A,-2A\}$ have density $O(1/p)$ and are not detected; Shor's algorithm tolerates them because it only needs the addition to be right on \emph{most} inputs. The windowed variant additionally has to handle the table row $A=\mathcal O$ (\cref{sec:windowed}).
\end{warnbox}

\subsection{The \ecdsafail{} challenge: harness, metric, validity}\label{sec:challenge}

\ecdsafail{} \cite{ecdsafail26} is a public benchmark launched by Eigen Labs on 30~May 2026: a Rust harness (adapted from the code Google released with \cite{babbush26}, Zenodo record \mbox{19597130}) in which contestants edit \code{src/point\_add/} to emit an \emph{op stream} for one mixed point addition, and a trusted evaluator (\code{eval\_circuit.rs}, which never imports contestant code) runs that stream on $9{,}024$ test inputs. The interface is: two $256$-qubit registers \code{target\_x}, \code{target\_y} holding the accumulator, two $256$-\emph{bit} classical registers \code{offset\_x}, \code{offset\_y} holding the addend, and the circuit must overwrite the targets with the affine sum. The gate set is $X,Z,\mathrm{CX},\mathrm{CZ},\mathrm{CCX},\mathrm{CCZ},\mathrm{SWAP}$, reset $R$, and $\Hmr$ (``Hadamard--measure--reset'': a demolition measurement in the $X$ basis whose outcome becomes a classical bit), with gates optionally conditioned on classical bits (\code{PushCondition}/\code{PopCondition}).

\paragraph{The score.} $S=\QT$ with
\begin{align*}
Q &= 1+\max\{\text{qubit index referenced in the op stream}\}\ \ (=\text{the high-water mark of live qubits}),\\
T &= \text{average over the $9{,}024$ shots of the number of CCX/CCZ gates \emph{executed}} .
\end{align*}
Executed, not emitted: a Toffoli under a classical condition counts only on the shots where the condition holds, which is what makes measurement-based uncomputation ``free'' (\cref{sec:mbu}) and why $T$ is not a static gate count, a $T$-depth, or a fault-tolerant hardware cost. Clifford gates, measurements and resets cost nothing.

\paragraph{The three channels.} Every run reports, for all $9{,}024$ inputs, (i) classical mismatches \code{cls} (wrong output point, per shot), (ii) phase-garbage batches \code{pha} (a non-zero accumulated global phase across the $64$-shot batch---a leftover $(-1)^{m\cdot f}$ from a measurement whose fix-up did not match, or a reset of a qubit that was not $\ket0$, which the simulator models by XOR-ing a random bit into the phase), and (iii) ancilla-garbage batches \code{anc} (a scratch qubit left non-zero at the end). A valid run is $0/0/0$; a valid run that beats the incumbent is \emph{promoted} and becomes a commit on \code{main}.

\paragraph{Fiat--Shamir islands, the tail nonce and $\lambda$.} The test inputs are not fixed: a SHAKE256 stream seeded with the domain tag, the op count and every field of every op draws the scalar pairs $(k_1,k_2)$, the points $[k_1]G,[k_2]G$, and also the random outcomes of every $R$ and $\Hmr$. A circuit therefore cannot be tuned to a known test set---but it can be tuned to \emph{some} test set. Every op stream ends with $48$ pairs of self-cancelling $X;X$ gates whose targets encode a $48$-bit \emph{tail nonce}: it costs nothing and changes the hash, hence the ``island'' (inputs plus measurement outcomes) on which the circuit is judged. A deliberately approximate circuit is wrong on a small set of \emph{hard inputs}; if $\lambda$ is the expected number of failing shots on a random island, a random nonce yields a clean island with probability roughly $e^{-\lambda}$, and contestants \emph{hunt} nonces (classical pre-filters on GPUs replay the circuit's own failure predicate at tens of thousands of nonces per second, then a full simulation confirms the survivors). The practitioners price every truncation in ``Toffoli per unit of $\lambda$'' and work against explicit $\lambda$ caps (18 and 19.5 in the September--October packages) \cite{ecdsafail-wiki}. This is the central integrity caveat of the whole leaderboard, stated by the paper itself: \emph{validator-clean is not value-exact}; the paper re-evaluates $42$ accepted submissions on a common $50{,}000$-input corpus and reports a retry-adjusted score $\QT/\hat p$, with $\hat p\approx0.998$ for the best circuit \cite{ecdsafail26}.

\begin{srcbox}[title={Where the material in Part~I comes from}]
The main paper \cite{ecdsafail26} (arXiv:2609.09582, v2 of 19~Sept.\ 2026, data cutoff 26~July 2026 09:21:55~UTC); its Supplemental Notes A--F, published as TeX and PDF in the GitHub repository \code{jieyilong/ecdsafail-supplemental-materials} \cite{ecdsafail-supp} (A: trajectory, B: anatomy of the optimisations, C: census of the $404$ scored commits, D: the best $\QT$ circuit \code{8e9c9a2}, E: the low-$Q$ circuit \code{b6f2b0a}, F: the window-compatible variant); the challenge repository \code{Layr-Labs/ecdsafail-challenge} \cite{ecdsafail-repo} (cloned at commit \code{c672dda}, 11~Oct.\ 2026); and the per-commit wiki \code{benhuang2025/ecdsafail-wiki} \cite{ecdsafail-wiki}, which documents all $584$ accepted commits through 7~Oct.\ 2026 with their diffs, $18$ concept pages, and leaderboard snapshots. Where the wiki's audit disagrees with the paper's census it says so, and so do we.
\end{srcbox}
